\chapter{进程管理}

\section{创建进程}
\begin{lstlisting}[style=demostyle,numbers=none]
import subprocess

subprocess.Popen('calc.exe')    # 打开计算器
subprocess.Popen('notepad.exe') # 打开记事本
subprocess.Popen('cmd.exe')     # 打开命令行
subprocess.Popen('resmon.exe')  # 打开资源管理器
subprocess.Popen('taskmgr.exe') # 打开任务管理器
subprocess.Popen('D:\\application\\wps\\WPS Office\\ksolaunch.exe') # 打开WPS
subprocess.Popen('D:/application/wps/WPS Office/ksolaunch.exe') # 打开WPS
\end{lstlisting}

\section{获取进程}
\subsection{显示所有进程及对应的名称}
\begin{lstlisting}[style=demostyle,numbers=none]
import psutil

# 获取所有进程的 PID
pids = psutil.pids()

# 遍历每个 PID，获取进程名称并打印
for pid in pids:
    try:
        # 实例化一个 Process 对象
        process = psutil.Process(pid)
        # 获取并打印进程名称
        print(f"PID: {pid}, Name: {process.name()}")
    except (psutil.NoSuchProcess, psutil.AccessDenied, psutil.ZombieProcess):
        # 忽略访问被拒绝、进程不存在或僵尸进程的情况
        pass
\end{lstlisting}
\begin{lstlisting}[style=demostyle,numbers=none]
PID: 0, Name: System Idle Process
PID: 4, Name: System
...
PID: 25836, Name: ChsIME.exe
PID: 25932, Name: chrome.exe
PID: 26128, Name: Code.exe
PID: 26236, Name: wsl.exe
PID: 26436, Name: conhost.exe
PID: 26608, Name: wps.exe
PID: 26612, Name: wslhost.exe
\end{lstlisting}

\subsection{根据名称获取进程ID}
\begin{lstlisting}[style=demostyle,numbers=none]
import psutil

def get_pids_by_name(process_name):
    pids = []
    for proc in psutil.process_iter(['pid', 'name']):
        try:
            if proc.info['name'] == process_name:
                pids.append(proc.info['pid'])
        except (psutil.NoSuchProcess, psutil.AccessDenied, psutil.ZombieProcess):
            pass  # 忽略访问被拒绝或进程不存在的情况
    return pids

process_name_to_find = "python"  # 替换为你想要查找的进程名称
pids = get_pids_by_name(process_name_to_find)

print(f"进程名称 '{process_name_to_find}' 的进程 ID 列表: {pids}")
\end{lstlisting}

\begin{lstlisting}[style=demostyle,numbers=none]
进程名称 'Code.exe' 的进程 ID 列表: [2108, 2692, 4676, 9472, 11052, 11964, 12184, 14440, 14872, 15052, 15584, 16896, 17392, 17784, 18924, 19040, 19132, 20708, 22544, 23080, 23868, 26128]
\end{lstlisting}


\begin{lstlisting}[style=demostyle,numbers=none]
import psutil

>>> psutil.pids() # 所有进程ID
[3865, 3864, 3863, 3856, 3855, 3853, 3776, ..., 45, 44, 1, 0]
>>> p = psutil.Process(3776) # 获取指定进程ID=3776，其实就是当前Python交互环境
>>> p.name() # 进程名称
'python3.6'
>>> p.exe() # 进程exe路径
'/Users/michael/anaconda3/bin/python3.6'
>>> p.cwd() # 进程工作目录
'/Users/michael'
>>> p.cmdline() # 进程启动的命令行
['python3']
>>> p.ppid() # 父进程ID
3765
>>> p.parent() # 父进程
<psutil.Process(pid=3765, name='bash') at 4503144040>
>>> p.children() # 子进程列表
[]
>>> p.status() # 进程状态
'running'
>>> p.username() # 进程用户名
'michael'
>>> p.create_time() # 进程创建时间
1511052731.120333
>>> p.terminal() # 进程终端
'/dev/ttys002'
>>> p.cpu_times() # 进程使用的CPU时间
pcputimes(user=0.081150144, system=0.053269812, children_user=0.0, children_system=0.0)
>>> p.memory_info() # 进程使用的内存
pmem(rss=8310784, vms=2481725440, pfaults=3207, pageins=18)
>>> p.open_files() # 进程打开的文件
[]
>>> p.connections() # 进程相关网络连接
[]
>>> p.num_threads() # 进程的线程数量
1
>>> p.threads() # 所有线程信息
[pthread(id=1, user_time=0.090318, system_time=0.062736)]
>>> p.environ() # 进程环境变量
{'SHELL': '/bin/bash', 'PATH': '/usr/local/bin:/usr/bin:/bin:/usr/sbin:/sbin:...', 'PWD': '/Users/michael', 'LANG': 'zh_CN.UTF-8', ...}
>>> p.terminate() # 结束进程
Terminated: 15 <-- 自己把自己结束了
\end{lstlisting}

\section{结束进程}
\begin{lstlisting}[style=demostyle,numbers=none]
import psutil

# 获取指定进程ID的进程对象
p = psutil.Process(pid)

# 结束进程
p.terminate()
\end{lstlisting}